\begin{solucion}{Contrato con \(\alpha\)}
Definimos:
\[
\begin{array}{c|l}
\toprule
\text{Simbolo} & \text{Significado} \\
\midrule
\alpha\in(0,1) & \text{fraccion asignada al minorista en el contrato.}\\
w\in\mathbb{R}_+ & \text{precio mayorista pagado por unidad al proveedor.}\\
q\in\mathbb{R}_+ & \text{cantidad elegida por el minorista.}\\
R(q)=p(q)q & \text{ingreso total generado por vender } q.\\
\bottomrule
\end{array}
\]
La variable de decision del minorista sigue siendo \(q\in\mathbb{R}_+\). El contrato coordina si la decision privada del minorista coincide con la decision eficiente de la cadena centralizada.

La utilidad del minorista bajo reparto es:
\[
\Pi_r(q)=\alpha R(q)-wq.
\]
Su condicion de primer orden es:
\[
\alpha R'(q)-w=0
\quad\Rightarrow\quad
R'(q)=\frac{w}{\alpha}.
\]
La cadena integrada elige la cantidad eficiente cuando:
\[
R'(q)=c.
\]
Por lo tanto, para coordinar:
\[
\frac{w}{\alpha}=c
\quad\Rightarrow\quad
\boxed{w=\alpha c}.
\]

En el caso base:
\[
c=1,
\qquad
\Pi_{SC}^C=2450.25.
\]
Si \(\alpha=0.35\), entonces:
\[
w=\alpha c=0.35(1)=0.35.
\]
La utilidad del minorista es:
\[
\Pi_r^A=\alpha\Pi_{SC}^C
=0.35(2450.25)
=857.5875.
\]
La utilidad del proveedor es:
\[
\Pi_s^A=(1-\alpha)\Pi_{SC}^C
=0.65(2450.25)
=1592.6625.
\]

Como en el caso descentralizado:
\[
\Pi_r^D=612.5625,
\qquad
\Pi_s^D=1225.125,
\]
ambos mejoran:
\[
857.5875>612.5625,
\qquad
1592.6625>1225.125.
\]

Para encontrar el rango aceptable:
\[
\alpha(2450.25)>612.5625
\quad\Rightarrow\quad
\alpha>0.25.
\]
Para el proveedor:
\[
(1-\alpha)(2450.25)>1225.125
\quad\Rightarrow\quad
1-\alpha>0.5
\quad\Rightarrow\quad
\alpha<0.5.
\]
Por lo tanto:
\[
\boxed{0.25<\alpha<0.50.}
\]
Cualquier valor dentro de ese rango coordina la cadena y deja a ambos agentes mejor que en el escenario descentralizado.
\end{solucion}
